\chapter{How to Use This Book}

\section{About the Wolf Prize in Mathematics}

The Wolf Prize in Mathematics was established by Ricardo Wolf in 1978 and is awarded annually (with occasional gaps) by the Wolf Foundation in Israel. Alongside the Fields Medal (awarded only to mathematicians under 40) and the Abel Prize, it is among the most prestigious international awards in mathematics. Unlike the Fields Medal, the Wolf Prize recognizes lifetime achievement, and its laureates span the full range of pure and applied mathematics from the late 19th century through the present day.

This book surveys the work of all 68 Wolf Prize laureates from 1978 through 2024. Each chapter focuses on one laureate, describing their main contributions, the historical context, the key mathematical ideas, and the lasting impact of their work.

\section{Mathematical Fields Represented}

The Wolf Prize laureates span virtually all areas of mathematics. The major fields, and the laureates associated with each, are:

\begin{description}
    \item[Algebra and Number Theory:] Siegel, Weil, Zariski, Erd\H{o}s, Langlands, Wiles, Tate, Margulis, Sarnak, Arthur, Deligne, Beilinson, Drinfeld, Lusztig
    \item[Algebraic Geometry:] Weil, Zariski, Kodaira, Mumford, Griffiths, Artin, Deligne
    \item[Analysis and PDEs:] Gelfand, Calder\'on, H\"ormander, Stein, Fefferman, Caffarelli, De Giorgi
    \item[Complex Analysis:] Ahlfors, Cartan, Carleson
    \item[Combinatorics and Computer Science:] Erd\H{o}s, Lov\'asz, Alon, Shamir
    \item[Differential Geometry and Topology:] Whitney, Chern, Milnor, Bott, Novikov, Donaldson, Eliashberg, Gromov
    \item[Dynamical Systems:] Kolmogorov, Sinai, Arnold, Moser, Furstenberg
    \item[Functional Analysis and Operator Theory:] Gelfand, Krein, Lax, Calder\'on
    \item[Geometry:] Gromov, Yau, Schoen, Mostow, Sullivan
    \item[Group Theory and Representation:] Thompson, Tits, Aschbacher, Lusztig, Margulis
    \item[Logic and Foundations:] G\"odel, Cohen, Shelah
    \item[Mathematical Physics:] Gelfand, Sinai, It\^o, Keller, Wiles
    \item[Probability:] Kolmogorov, It\^o, Le Gall, Lawler
    \item[Topology:] Eilenberg, Hirzebruch, Milnor, Smale, Novikov, Bott, Serre, Sullivan
    \item[Wavelets and Signal Processing:] Daubechies
\end{description}

\section{How to Read This Book}

\subsection{Mathematical Maturity Levels}

Each chapter is labelled with a suggested background level:

\begin{itemize}
    \item \textbf{[B] Basic:} Assumes only standard undergraduate mathematics (calculus, linear algebra, basic abstract algebra). Suitable for advanced undergraduates.
    \item \textbf{[I] Intermediate:} Assumes one or two upper-division or introductory graduate courses (e.g., real analysis, algebraic topology, complex analysis). Suitable for beginning graduate students or advanced undergraduates with a targeted background.
    \item \textbf{[A] Advanced:} Assumes graduate-level familiarity with the relevant field. Suitable for graduate students and researchers.
\end{itemize}

Most chapters contain material at multiple levels; the label reflects the minimum background needed to follow the main exposition.

\subsection{Suggested Reading Paths}

Rather than reading linearly, undergraduates may want to follow one of these paths:

\textbf{Path 1 --- Number Theory and Algebra:}
Weil $\to$ Tate $\to$ Deligne $\to$ Langlands $\to$ Arthur $\to$ Wiles $\to$ Sarnak

\textbf{Path 2 --- Analysis and PDEs:}
Gelfand $\to$ Krein $\to$ Calder\'on $\to$ H\"ormander $\to$ Stein $\to$ Fefferman $\to$ Caffarelli

\textbf{Path 3 --- Geometry and Topology:}
Whitney $\to$ Milnor $\to$ Bott $\to$ Gromov $\to$ Yau $\to$ Donaldson $\to$ Eliashberg

\textbf{Path 4 --- Probability and Applications:}
Kolmogorov $\to$ It\^o $\to$ Sinai $\to$ Le Gall $\to$ Lawler $\to$ Daubechies

\textbf{Path 5 --- Combinatorics and Computer Science:}
Erd\H{o}s $\to$ Lov\'asz $\to$ Alon $\to$ Shamir $\to$ Daubechies

\textbf{Path 6 --- Groups and Symmetry:}
Thompson $\to$ Tits $\to$ Aschbacher $\to$ Lusztig $\to$ Margulis

\subsection{Structure of Each Chapter}

Every chapter is organized into the following sections (not all chapters include every section):

\begin{enumerate}
    \item \textbf{Prerequisites:} The background knowledge assumed for that chapter.
    \item \textbf{Main Contributions:} A high-level overview of the laureate's work.
    \item \textbf{Origin of the Problem:} The historical and mathematical context.
    \item \textbf{How It Was Solved:} The key ideas, insights, and techniques.
    \item \textbf{Mathematical Theory:} Definitions, theorems, and proofs.
    \item \textbf{Impact:} How the work influenced mathematics, science, and technology.
    \item \textbf{Frequently Asked Questions:} Common questions answered in plain language.
    \item \textbf{Exercises:} Practice problems labelled by difficulty.
    \item \textbf{Timeline:} Key events in the laureate's life and career.
    \item \textbf{Glossary:} Definitions of specialized terms used in the chapter.
    \item \textbf{Citations:} Primary sources and further reading.
\end{enumerate}

\section{Cross-Chapter Dependencies}

Some chapters build on ideas introduced in others. The dependency graph below shows which chapters should be read before others for maximum comprehension.

\begin{center}
\begin{tabularx}{\linewidth}{|l|X|}
\hline
\textbf{Read First} & \textbf{Then Read} \\ \hline
Gelfand (ch01) & Krein (ch10), Lax (ch18), Calder\'on (ch21) \\ \hline
Weil (ch04) & Deligne (ch46), Langlands (ch30), Wiles (ch31) \\ \hline
Cartan (ch05) & Serre (ch37), Leray (ch03) \\ \hline
Kolmogorov (ch06) & Sinai (ch33), It\^o (ch17), Arnold (ch38) \\ \hline
Ahlfors (ch07) & Kodaira (ch13), Yau (ch50) \\ \hline
Zariski (ch08) & Mumford (ch48), Artin (ch54) \\ \hline
Eilenberg (ch15) & Serre (ch37), Bott (ch36), Milnor (ch22) \\ \hline
Selberg (ch16) & Langlands (ch30), Arthur (ch56), Sarnak (ch55) \\ \hline
Hirzebruch (ch19) & Bott (ch36), Donaldson (ch63) \\ \hline
Thompson (ch26) & Aschbacher (ch51), Tits (ch28) \\ \hline
Moser (ch29) & Arnold (ch38), Sinai (ch33) \\ \hline
It\^o (ch17) & Le Gall (ch61), Lawler (ch62) \\ \hline
Milnor (ch22) & Smale (ch44), Donaldson (ch63), Eliashberg (ch64) \\ \hline
\end{tabularx}
\end{center}

\section{Undergraduate-Friendly Entry Points}

If you are an undergraduate looking for accessible entry points into this book, we recommend starting with these chapters. They require minimal specialized background and give a feel for the breadth of the Wolf Prize:

\begin{enumerate}
    \item \textbf{Paul Erd\H{o}s (ch12):} Combinatorics, graph theory, probabilistic method. Requires only basic discrete math and probability.
    \item \textbf{Noga Alon (ch67):} Combinatorial Nullstellensatz, expander graphs. Requires linear algebra and basic combinatorics.
    \item \textbf{Adi Shamir (ch68):} RSA cryptography, secret sharing. Requires elementary number theory.
    \item \textbf{Ingrid Daubechies (ch66):} Wavelets, signal processing. Requires Fourier analysis and linear algebra.
    \item \textbf{Kiyoshi It\^o (ch17):} Stochastic calculus. Requires calculus and basic probability.
    \item \textbf{L\'aszl\'o Lov\'asz (ch34):} Combinatorics, graph theory, theoretical computer science. Requires discrete math.
    \item \textbf{John Milnor (ch22):} Topology and geometry. Requires basic topology and real analysis.
    \item \textbf{Peter Lax (ch18):} Analysis, PDEs, applied math. Requires calculus and real analysis.
    \item \textbf{Yakov Sinai (ch33):} Dynamical systems, ergodic theory. Requires real analysis and basic probability.
    \item \textbf{Stephen Smale (ch44):} Topology, dynamical systems. Requires basic topology and real analysis.
\end{enumerate}

Once you have read a few of these, you will have the background to approach the Intermediate-level chapters and, from there, the Advanced ones.

\section{A Note on Mathematical Rigor}

This book aims to be mathematically precise but not encyclopedic. Definitions and theorems are stated carefully, and key proofs are sketched or given in full where they illustrate an important technique. However, the goal is understanding, not formal completeness. Readers who want to see the full technical details should consult the primary sources listed in each chapter's Citations.
